%%%%%%%%%%%%%%%%%%%%%%%%%%%%%%%%%%%%%%%%%%%%%%%%%%%%%%%%%%%%%%
%%%%%%%%%%%%%%%%%%%%%% MA-0711 Lógica %%%%%%%%%%%%%%%%%%%%%%%%
%%%%%%%%%%%%%%%%%%%%%%%%%%%%%%%%%%%%%%%%%%%%%%%%%%%%%%%%%%%%%%
\newpage
\subsection*{MA-0711 Lógica}
\addcontentsline{toc}{subsection}{MA-0711 Lógica}

\hypertarget{MA0711}{}
\begin{refsection}
\begin{table}[H]
\centering{
\begin{tabular}{|ll|}
\hline
\textbf{Año:} IV  & \textbf{Requisitos:} MA-0451 Principios de análisis en una variable; \\ & MA-0461 Álgebra Lineal II \\ 
\textbf{Ciclo:} Optativa  & \textbf{Correquisitos:} Ninguno \\ 
\textbf{Tipo de curso:} Teórico & \textbf{Horas presenciales por semana:} 5 \\ 
\textbf{Créditos:} 5 & \textbf{Horas de trabajo independiente por semana:} 10 \\ \hline
\end{tabular}
}
\end{table}

\subsubsection*{Descripción del curso}

La lógica matemática es una de las áreas fundamentales de la matemática. Su objetivo es estudiar los lenguajes y sistemas formales utilizados para expresar razonamientos matemáticos, así como las estructuras que los interpretan. La interacción entre la sintaxis y la semántica permite analizar las consecuencias lógicas de un conjunto de axiomas y comprender cómo estos determinan el comportamiento de las estructuras que los satisfacen.

Por ejemplo, la teoría de grupos estudia las propiedades que se deducen de los axiomas que definen un grupo. De manera análoga, los modelos formales de computación permiten investigar qué problemas pueden resolverse mediante algoritmos y cuáles son las limitaciones inherentes de la computación. De este modo, la lógica matemática mantiene una estrecha relación con numerosas áreas de la matemática y de las ciencias de la computación.

Este curso está orientado al estudio de los fundamentos de la disciplina. Se introducirán las nociones básicas de teoría de conjuntos, incluyendo ordinales y cardinales, así como la sintaxis y la semántica del cálculo de predicados de primer orden. Estas herramientas permitirán demostrar los teoremas de completitud y de compacidad, resultados centrales de la lógica matemática.

Finalmente, se introducirán nociones básicas de teoría de modelos: teoría, teoría completa, modelo, isomorfismo y extensión elemental, culminando con los teoremas de Löwenheim--Skolem y el criterio de Tarski--Vaught.

\subsubsection*{Objetivos}

\textbf{General:} Desarrollar los conocimientos y herramientas fundamentales de la lógica matemática para analizar y formalizar razonamientos matemáticos y demostrar los principales resultados del cálculo de predicados de primer orden.

\textbf{Específicos:} Al finalizar el curso se espera que la persona estudiante sea capaz de:

\begin{enumerate}
\item Realizar aritmética básica de ordinales y cardinales.
\item Entender el concepto formal de lenguaje y fórmula de primer orden, la forma en que se construyen y utilizar esto para demostrar resultados sobre la sintaxis de un lenguaje.
\item Entender el concepto formal de estructura de primer orden, así como su relación semántica con una fórmula.
\item Relacionar los conceptos de sintaxis y semántica mediante el teorema de completitud.
\item Aplicar el teorema de compacidad a casos concretos de estructuras de primer orden naturales.
\end{enumerate}

\subsubsection*{Contenidos}

\begin{enumerate}

\item \textbf{Ordinales y cardinales (3 semanas):}
\begin{enumerate}
    \item Conjuntos, relaciones y funciones; equipotencia y enumerabilidad.
    \item Construcción de ordinales y cardinales; aritmética ordinal y cardinal elemental.
\end{enumerate}

\item \textbf{Lógica de primer orden (4 semanas):}
\begin{enumerate}
    \item Lenguaje del cálculo de predicados; estructuras e interpretaciones.
    \item Sintaxis y semántica: términos, fórmulas, satisfacción y consecuencia lógica.
    \item Sistema formal de demostración; formalización de conceptos matemáticos.
\end{enumerate}

\item \textbf{Teorema de completitud y teorema de compacidad (4 semanas):}
\begin{enumerate}
    \item Demostración del teorema de completitud de Gödel.
    \item Teorema de compacidad y consecuencias fundamentales.
    \item Puente entre sintaxis y semántica; aplicaciones introductorias.
\end{enumerate}

\item \textbf{Nociones básicas de teorías y estructuras (4 semanas):}
\begin{enumerate}
    \item Teoría, modelo, subestructura, isomorfismo y extensión elemental.
    \item Teoremas de Löwenheim--Skolem; criterio de Tarski--Vaught.
    \item Aplicaciones en álgebra y otras áreas de la matemática.
\end{enumerate}

\end{enumerate}

\subsubsection*{Metodología}

\noindent \textbf{Lineamientos metodológicos:} La persona docente a cargo de este curso debe respetar las siguientes pautas:

\begin{enumerate}
    \item Desarrollar el curso principalmente mediante \textbf{clases magistrales presenciales}, con participación activa del estudiantado.
    \item Complementar con \textbf{sesiones de ejercicios} para consolidar conceptos teóricos y fortalecer habilidades de demostración.
    \item Pueden abordarse temas opcionales (introducción axiomática a teoría de conjuntos, recursión computable) mediante talleres o proyectos, según el tiempo disponible.
    \item Promover evaluación formativa y sumativa coherente con el marco pedagógico de la carrera; se sugiere incluir al menos dos exámenes parciales.
    \item Cuando las autoridades sanitarias o institucionales impongan restricciones, adaptar las actividades según normativa vigente.
\end{enumerate}

\textbf{Sugerencias metodológicas:} Adicionalmente, se tienen las siguientes recomendaciones para la persona docente a cargo de este curso:

\begin{enumerate}
    \item Fomentar la lectura de textos de referencia como \cite{HiL19}, \cite{EFT94} y \cite{CL2000}.
    \item Utilizar plataformas institucionales de mediación virtual para material y entregas, cuando corresponda.
    \item Orientar proyectos individuales o grupales sobre temas opcionales del programa (teoría de conjuntos axiomática, computabilidad).
\end{enumerate}

\nocite{HiL19}
\nocite{CL2000}
\nocite{EFT94}
\nocite{Mar02}
\nocite{Poi00}
\nocite{Kun80}

\printbibliography[heading=subbibliography,section=\the\value{refsection}]
\end{refsection}
